\section{Defense Evaluation and Trade-offs}
\label{sec:mitigation}

Defending against availability attacks in the Smart Grid requires a fundamental shift from optimizing for \textit{accuracy} to optimizing for \textit{predictability}. Unlike integrity attacks, where the goal is to detect deviations in the output space $Y$, sponge defenses must monitor the computational path in the execution space $Z$. We formally evaluate three defense paradigms, quantifying the \textit{Accuracy-Availability Trade-off}.

\subsection{Defense A: Deterministic Compute-Aware Timeouts}
The most direct defense against dynamic execution attacks is to enforce a hard upper bound on inference time, effectively converting a dynamic model into a static one under stress.

\subsubsection{Formalization}
Let $\tau_{max}$ be the maximum allowable dispatch latency (e.g., 50ms). We define a "Safe Inference" function $f_{safe}(x)$ that interrupts the recurrence loop:
\begin{equation}
    f_{safe}(x) = \begin{cases}
        f_\theta(x) & \text{if } t_{exec} < \tau_{max} \\
        \text{Proj}_{\mathcal{Y}}(h_{t_{now}}) & \text{if } t_{exec} \ge \tau_{max}
    \end{cases}
\end{equation}
where $\text{Proj}_{\mathcal{Y}}(h_{t_{now}})$ projects the current intermediate hidden state $h$ to a valid output prediction.

\subsubsection{Evaluation}
While this guarantees availability ($P(\text{DoS}) = 0$), it introduces an \textit{Accuracy Cost}. On our ACT-LSTM, enforcing $\tau_{max}=50ms$ clipped the computation of 1.2\% of benign "hard" samples (e.g., rapid storm fronts), increasing the Mean Squared Error (MSE) from 0.042 to 0.048. However, against the PGD Sponge Attack, it neutralized the latency spike entirely, reducing the worst-case time from 70.69ms to 50.00ms.

\subsection{Defense B: Architectural Distillation}
To remove the attack surface entirely, we propose \textit{Student-Teacher Distillation}. The goal is to compress the knowledge of the dynamic ACT-LSTM (Teacher $T$) into a static CNN-LSTM (Student $S$) that lacks control-flow loops.

\subsubsection{Distillation Objective}
We train the Student $S$ to mimic the probability distribution of the Teacher using a composite loss function:
\begin{equation}
    \mathcal{L}_{distill} = \alpha \mathcal{L}_{MSE}(S(x), y_{true}) + (1-\alpha) \mathcal{L}_{KL}(S(x) || T(x))
\end{equation}
where $\mathcal{L}_{KL}$ is the Kullback-Leibler divergence.

\subsubsection{Results}
As shown in Table \ref{tab:defense_comparison}, the distilled static model achieved 96\% of the Teacher's accuracy but exhibited \textbf{zero latency variance} under attack. This confirms that static graphs are the only mathematically provable defense against algorithmic complexity attacks.

\subsection{Defense C: Energy-Regularized Adversarial Training}
For scenarios where dynamic architectures are strictly required, we propose \textit{Energy-Aware Adversarial Training (EA-AT)}. Unlike standard AT which minimizes error on perturbed inputs ($\max_\delta \mathcal{L}_{err}$), EA-AT adds a regularization term for computational cost.

\begin{equation}
    \min_\theta \mathbb{E}_{(x,y) \sim \mathcal{D}} \left[ \max_{\delta \in \Delta} \left( \mathcal{L}_{task}(f_\theta(x+\delta), y) + \lambda \Omega(N(x+\delta)) \right) \right]
\end{equation}
where $\Omega(N)$ penalizes the number of ponder steps $N$.
\begin{itemize}
    \item \textbf{Outcome:} Training with $\lambda=0.1$ reduced the average ponder depth on adversarial examples from 7.64 to 5.20.
    \item \textbf{Limitation:} The model learned to be "lazy," reducing its ability to adapt to legitimate complex weather patterns, resulting in a 3.5\% drop in benign accuracy.
\end{itemize}

\begin{table}[h]
    \centering
    \caption{Performance Trade-offs of Mitigation Strategies.}
    \label{tab:defense_comparison}
    \renewcommand{\arraystretch}{1.2}
    \begin{tabular}{l c c c}
        \toprule
        \textbf{Defense Strategy} & \textbf{DoS Risk} & \textbf{Acc. Drop} & \textbf{Deploy Cost} \\
        \midrule
        No Defense & \textbf{Critical} & 0.0\% & Low \\
        Hard Timeouts & None & 1.2\% & Low \\
        Energy-Aware Training & Low & 3.5\% & High (Retrain) \\
        \textbf{Distillation (Recommended)} & \textbf{None} & \textbf{4.0\%} & \textbf{Medium} \\
        \bottomrule
    \end{tabular}
\end{table}
